\documentclass[10pt,a4paper]{amsart}
\usepackage[margin=19mm]{geometry}
\usepackage{amsmath,amssymb,mathtools}
\usepackage[scr=rsfs,cal=euler]{mathalfa}
\usepackage{xfrac}
\usepackage[unicode,hidelinks,bookmarks=false]{hyperref}
\newcommand{\ie}{i.e.\ }
\newcommand{\cf}{cf.\ }
\newcommand{\resp}{respectively\ }
\newcommand{\Subsection}{\subsection}
\providecommand{\nobreakdash}{\nobreak\hbox{-}}
\setlength{\emergencystretch}{2em}
\multlinegap0pt
\usepackage{bm}
\usepackage{bbm}
\usepackage{dsfont}
\usepackage{xspace}
\usepackage{cancel}
\usepackage{mathtools}
\usepackage[shortlabels]{enumitem}
\usepackage{amsthm}
\makeatletter
\@ifundefined{theorem}{\newtheorem{theorem}{Theorem}}{}
\makeatother
\DeclareMathOperator{\coker}{coker}                   % Cokernel  
\DeclareMathOperator{\rank}{rank}
\newcommand{\Ring}{R}                          % Discrete valuation ring
\newcommand{\coboundary}{d}                           % Coboundary operator
\newcommand{\sheaf}[1]{\mathcal{#1}}                  % Sheaf symbol
\newcommand{\unif}{\pi}                        % Uniformizer
\newcommand{\valpi}{\mathrm{val}_{\unif}}                   % Valuation with respect to uniformizer
\title[Precision-Graded Cohomology and Arithmetic Persistence for N]{Precision-Graded Cohomology and Arithmetic Persistence for Network Sheaves}
\author{Robert Ghrist, Cassie Ding}
\date{Source: arXiv:2511.00677v1 (CC BY 4.0)}
\begin{document}
\maketitle

\noindent\emph{Source and licence.} Precision-Graded Cohomology and Arithmetic Persistence for Network Sheaves. Version arXiv:2511.00677v1 (CC BY 4.0), distributed under the Creative Commons Attribution 4.0 International licence, as recorded in the arXiv licence metadata for this exact version and shown on the abstract page https://arxiv.org/abs/2511.00677v1. Full rights notice: licenses/arxiv-2511.00677-CC-BY-4.0.txt.\par\smallskip

\noindent\textit{Editorial scope.} Counted unit: the Theorem whose source label is \texttt{thm:cycle-barcode} in the article (source0.tex lines 1413--1423, source label \texttt{thm:cycle-barcode}), delivered together with the article's own complete original proof of it (source0.tex lines 1425--1449, one \texttt{proof} environment of 25 lines). No mathematics is written, repaired or re-derived: statement and proof are the article's own text line for line. Required context retained uncounted: none. Every definition, display and label of those lines is retained unchanged. Omitted from this package and not redistributed: 1--1412, 1424--1424, 1450--1798. Nothing inside a retained excerpt is omitted or altered: every retained body is delivered exactly as the article's lines are written, so no omission receipt of any kind accompanies this package. Citations: the retained excerpts carry no citation command at all, so no bibliography entry of the article is transcribed and no reference is redistributed. Reference adaptation: none was needed and none was made; every reference inside the delivered unit resolves inside it, so no unresolved reference or citation is produced. Printed slips kept literal: no printed word, symbol, hypothesis or number of the counted statement or of its proof was altered. Layout and class: the article's own class is \texttt{amsart} with packages including babel, fontenc, inputenc, libertinus, parskip, geometry, amsmath, amsthm, amsfonts, amssymb, mathtools, bbm, enumerate, graphicx; the wrapper is the lane's pinned \texttt{amsart} editorial wrapper at 10pt on A4, which restates the theorem-like environments the retained text needs and adds the article's own macro definitions that the retained text needs (\texttt{\textbackslash Ring}, \texttt{\textbackslash coboundary}, \texttt{\textbackslash coker}, \texttt{\textbackslash rank}, \texttt{\textbackslash sheaf}, \texttt{\textbackslash unif}, \texttt{\textbackslash valpi}), with extra wrapper packages (bm, bbm, dsfont, xspace, cancel, mathtools, enumitem). The delivered numbering is the unit's own. Assets: no retained span contains a figure environment, a graphics-inclusion command, a PDF-page inclusion, a table or a TikZ picture; no image file is copied into this package, the package declares no asset dependency and no image file is redistributed with it. Licence: version arXiv:2511.00677v1 of this article is distributed under the Creative Commons Attribution 4.0 International licence, as recorded in the arXiv licence metadata for this exact version and shown on the abstract page https://arxiv.org/abs/2511.00677v1. A full rights notice is delivered at licenses/arxiv-2511.00677-CC-BY-4.0.txt. Characters: every retained excerpt is pure ASCII, so the delivered engine, XeLaTeX with its default Latin Modern text font, reports no missing character.
\begin{theorem}[Cycle Barcode]
\label{thm:cycle-barcode}
For the cycle graph $C_n$ with a rank-1 unit sheaf $\sheaf{F}$ having holonomy $h(C_n)$, the cohomology decomposes as follows:
\begin{enumerate}[label=\textup{(\roman*)}]
\item If $h(C_n) = 1$ (the constant sheaf): $H^0(C_n;\sheaf{F}) \cong \Ring$ and $H^1(C_n;\sheaf{F}) \cong \Ring$ (both free). The barcode consists of one infinite bar $[0,\infty)$ at degree $1$.

\item If $h(C_n) - 1$ is a unit in $\Ring$: $H^0(C_n;\sheaf{F}) = 0$ and $H^1(C_n;\sheaf{F}) = 0$. The barcode is empty.

\item If $\valpi(h(C_n)-1) = a > 0$: $H^0(C_n;\sheaf{F}) = 0$ and $H^1(C_n;\sheaf{F}) \cong \Ring/\unif^a$ (pure torsion). The barcode consists of a single finite bar $[0,a)$.
\end{enumerate}
\end{theorem}

\begin{proof}
Label vertices $v_0, \ldots, v_{n-1}$ and oriented edges $e_i = (v_i, v_{i+1})$ with indices modulo $n$. The coboundary $\coboundary: C^0 \to C^1$ is represented by the $n \times n$ matrix
\[
\coboundary = \begin{pmatrix}
-1 & m_0 & 0 & \cdots & 0 \\
0 & -1 & m_1 & \cdots & 0 \\
\vdots & \vdots & \ddots & \ddots & \vdots \\
0 & 0 & \cdots & -1 & m_{n-2} \\
m_{n-1} & 0 & \cdots & 0 & -1
\end{pmatrix}.
\]
Computing the determinant by cofactor expansion along the last row yields
\[
\det(\coboundary) = 1 - h(C_n) = u \cdot (h(C_n) - 1)
\]
for a unit $u \in \Ring^\times$. Consequently, $\valpi(\det(\coboundary)) = \valpi(h(C_n) - 1)$.

A cocycle $(x_0, \ldots, x_{n-1}) \in \ker \coboundary$ satisfies $m_i x_{i+1} = x_i$ for all $i$ (indices mod $n$), forcing $x_{i+1} = m_i^{-1}x_i$. Going around the cycle gives $x_0 = (m_0 \cdots m_{n-1})^{-1}x_0 = h(C_n)^{-1}x_0$. Thus $(h(C_n) - 1)x_0 = 0$ in $\Ring$.

\textit{Case (i): $h(C_n) = 1$.} Then $\det(\coboundary) = 0$, so $\coboundary$ is not invertible. The kernel is $\ker \coboundary = \{(x,x,\ldots,x) : x \in \Ring\} \cong \Ring$. Since $\rank(\coboundary) = n-1$ (as the $(n-1) \times (n-1)$ upper-left submatrix is invertible), the Smith form is $\mathrm{diag}(1,\ldots,1,0)$ with $(n-1)$ ones. Thus $H^1 = \coker(\coboundary) \cong \Ring$ is free.

\textit{Case (ii): $h(C_n) - 1 \in \Ring^\times$.} Then $\det(\coboundary)$ is a unit, so $\coboundary$ is invertible. Thus $\ker \coboundary = 0$ and $\coker(\coboundary) = 0$.

\textit{Case (iii): $\valpi(h(C_n)-1) = a > 0$.} Write $h(C_n) - 1 = u\unif^a$ with $u \in \Ring^\times$. Then $\det(\coboundary) = u\unif^a$ is neither zero nor a unit. The equation $(h(C_n) - 1)x_0 = 0$ has only the trivial solution $x_0 = 0$ in the free module $\Ring$, so $\ker \coboundary = 0$. The Smith normal form is $\mathrm{diag}(1,\ldots,1,\unif^a)$ (with $(n-1)$ ones), giving $H^1 \cong \Ring/\unif^a$, a single bar of length $a$.
\end{proof}

\end{document}
